% Generated draft: requires mathematical and source review.
\documentclass[11pt]{article}
\usepackage[margin=25mm]{geometry}
\usepackage{fontspec}
\setmainfont{Latin Modern Roman}
\usepackage{amsmath,amssymb,mathtools,unicode-math}
\setmathfont{Latin Modern Math}
\usepackage{enumitem,needspace,xurl,hyperref}
\setlength{\parindent}{0pt}
\setlength{\parskip}{5pt}
\setlength{\emergencystretch}{3em}
\newcommand{\N}{\mathbb{N}}
\newcommand{\Z}{\mathbb{Z}}
\newcommand{\Q}{\mathbb{Q}}
\newcommand{\R}{\mathbb{R}}
\newcommand{\C}{\mathbb{C}}
\newcommand{\F}{\mathbb{F}}
\newcommand{\A}{\mathbb{A}}
\newcommand{\PP}{\mathbb{P}}
\newcommand{\E}{\mathbb{E}}
\newcommand{\eps}{\varepsilon}
\newcommand{\dd}{\,\mathrm{d}}
\begin{document}
\section*{UnsolvedMath ID: 30002952}
\textbf{Source problem code:} catalogue entry
\subsection{Definitions and mathematical statement}
A branching random walk combines a genealogical branching process with spatial displacements along its genealogical edges. In the usual Galton--Watson formulation, criticality means that the offspring random variable \(\xi\) satisfies \(\E[\xi]=1\). This condition alone does not specify the offspring law, the displacement law, or the conditioning of the process. The supplied record refers to dimension \(d<6\), but does not specify the underlying spatial state space or the precise construction of the trace.

For reference, on a finite connected electrical network \(G=(V,E)\) with edge resistances \(r_e>0\), the effective resistance between distinct vertices \(a,b\in V\) is
\[
R_{\mathrm{eff}}^G(a,b)=\inf_{\theta:\,\operatorname{div}\theta=\delta_a-\delta_b}\sum_{e\in E}r_e\theta(e)^2,
\]
where each unoriented edge is counted once and \(\theta\) is an antisymmetric flow on oriented edges. Applying this definition to a branching-random-walk trace requires specifying how the trace becomes an electrical network, its terminals or boundary sets, and, for an infinite network, the boundary convention.

\textbf{Problem as supplied.} For the trace of critical branching random walk in dimension \(d<6\), determine the precise polynomial exponent governing the growth of its effective resistance \(R(n)\).

\textbf{Unresolved formulation.} The record does not define \(n\), the resistance terminals, the network conductances, or whether \(R(n)\) denotes a random resistance or an averaged quantity. Once these choices have been fixed, one possible meaning of a polynomial exponent would be a number \(\rho(d)\) satisfying
\[
\frac{\log R(n)}{\log n}\longrightarrow\rho(d)\qquad(n\to\infty),
\]
with a specified mode of convergence when \(R(n)\) is random. This is only a possible interpretation, not a formulation established by the supplied record. An exponent for \(\E[R(n)]\), almost-sure growth, convergence in probability, and two-sided power bounds are distinct targets. The supplied evidence does not determine which target is intended, so a unique precise conjecture cannot yet be stated.
\subsection{Short English statement}
The question asks how effective resistance grows in a critical branching-random-walk trace in dimension \(d<6\). The supplied record leaves the trace model, electrical network, meaning of \(n\), and probabilistic interpretation of the exponent unspecified. These choices must be recovered from the cited sources before a precise conjecture can be fixed.
\subsection{Sources}
\begin{itemize}
\item Local UnsolvedMath registry: permanent ID \(30002952\), title \emph{Resistance Exponents of Critical Branching Random-Walk Traces}, collection \emph{Oberwolfach Reports Open Problems}, exact source problem code \texttt{catalogue entry}. The supplied proposed year is 2015; no individual proposer is identified.
\item Supplied source citation: \emph{Scaling Limits in Models of Statistical Mechanics (2016)}, Oberwolfach Reports, DOI 10.4171/owr/2015/40, pp.\ 2329--2398. \url{https://doi.org/10.4171/owr/2015/40}. The literature-source list labels this same reference \emph{Oberwolfach Report 40, 2015}; the year discrepancy is unresolved.
\item Supplied literature reference: Nachmias--Járai, \emph{Electrical Resistance of the Low Dimensional Critical Branching Random Walk}. \url{https://arxiv.org/abs/1305.1092}.
\item Supplied literature reference: Járai--Nachmias, \emph{Logarithmic Correction to Resistance}. \url{https://arxiv.org/abs/2006.09005}.
\item Supplied dated literature curation, checked at \texttt{2026-08-21}: the record labels the problem open and reports sublinear upper bounds and work on critical corrections, while reporting no established precise polynomial growth exponent for \(d<6\). This assessment is supplied dated curation derived from the OWR material, not an independently verified literature assessment. The background explicitly requests independent verification of time-sensitive status and cited claims.
\end{itemize}
\end{document}
